% Architecture figure: replaces the whole figure* block (tikzpicture + caption)
% in the draft. Uses the \mdfmap/\mdslabs/\mdimstack macros and colours that are
% already in your preamble, and images/thumbs/{lr_rgb,sr_rgb}.png.
% Only change to the drawing: the sum node and the learned-knee box moved right
% (12.55 -> 12.75, 13.95 -> 14.45), so the arrow between them is visible.
% No blank lines inside the tikzpicture: with the linenumbers class option each
% one would end the build with "Improper \prevdepth".
\begin{figure*}
\centering
\begin{tikzpicture}
% ===================================================================== (a)
\def\hL{1.45}   % LR grid
\def\hH{2.35}   % SR grid
\def\ys{-3.15}  % skip row
\def\yi{-2.65}  % centre of the input image
\def\yo{-2.80}  % centre of the output image
\def\lrs{1.65}  % LR image side; the SR image is exactly twice this
\node[mdpanel] at (0,2.05) {(a)};
% input (lower left), Г turn up and into the stretch
\mdimstack{in}{0}{\yi}{\lrs}{lr_rgb}
\node[mdop,below=2pt of in-s,xshift=3pt] {LR, 4 bands};
\mdfmap{S}{2.70}{0}{0.72}{\hL}{mdLR}
\mdslabs{2.70}{0}{0.72}{\hL}{8}{mdLR}
\draw[mdflow,rounded corners=5pt] ($(in-s)+(0.105,\lrs+0.21)$) -- (0.5*\lrs+0.105,0) -- (S-w)
  node[mdop,pos=0.55,above=1pt] {asinh\\ 8 knees};
\node[mdch,above=1pt of S-n] {32};
% entry conv
\mdfmap{F}{4.65}{0}{0.72}{\hL}{mdLR}
\draw[mdflow] (S-e) -- (F-w) node[mdop,midway,above=1pt] {conv\\ $3\!\times\!3$};
\node[mdch,above=1pt of F-n] {32};
% residual trunk
\node[mdbox,draw=mdRes!75!black,fill=mdRes!35,minimum width=1.85cm,minimum height=\hL cm]
  (res) at (7.02,0) {residual\\ block $\times16$};
\draw[mdflow] (F-e) -- (res.west);
% tail conv + shuffle
\mdfmap{T}{8.80}{0}{0.59}{\hL}{mdLR}
\draw[mdflow] (res.east) -- (T-w) node[mdop,midway,above=1pt] {conv\\ $3\!\times\!3$};
\node[mdch,above=1pt of T-n] {16};
\mdfmap{M}{11.10}{0}{0.33}{\hH}{mdHR}
\draw[mdflow] (T-e) -- (M-w) node[mdop,midway,above=1pt] {pixel\\ shuf\/f\/le};
\node[mdch,above=1pt of M-n] {4};
% sum, knee, Г turn down into the output (exactly 2x the input image)
\node[mdsum] (sum) at (12.75,0) {$+$};
\draw[mdflow] (M-e) -- (sum);
\node[mdbox,draw=mdKnee!80!black,fill=mdKnee!18,minimum height=\hL cm] (knee) at (14.45,0)
  {learned knee\\[2pt] $k_b \sinh y$};
\draw[mdflow] (sum) -- (knee.west);
\mdimstack{out}{14.48}{\yo}{2*\lrs}{sr_rgb}
\draw[mdflow,rounded corners=5pt] (knee.east) -- (14.48+\lrs+0.105,0) -- ($(out-s)+(0.105,2*\lrs+0.21)$);
\node[mdop,below=2pt of out-s,xshift=3pt] {SR, 4 bands};
% per-band skip
\mdfmap{K}{4.65}{\ys}{0.46}{\hL}{mdLR}
\mdslabs{4.65}{\ys}{0.46}{\hL}{8}{mdLR}
\mdfmap{KC}{6.55}{\ys}{0.33}{\hL}{mdLR}
\mdfmap{KS}{8.65}{\ys}{0.07}{\hH}{mdHR}
\mdfmap{KA}{11.10}{\ys}{0.33}{\hH}{mdHR}
\draw[mdflow,rounded corners=4pt] ($(S-s)+(0.06,0)$) |- (K-w)
  node[mdop,pos=0.72,above=1pt] {band $b$};
\node[mdch,above=1pt of K-n] {8};
\draw[mdflow] (K-e) -- (KC-w) node[mdop,midway,above=1pt] {conv\\ $5\!\times\!5$};
\node[mdch,above=1pt of KC-n] {4};
\draw[mdflow] (KC-e) -- (KS-w) node[mdop,midway,above=1pt] {pixel\\ shuf\/f\/le};
\node[mdch,above=1pt of KS-n] {1};
\draw[mdflow] (KS-e) -- (KA-w) node[mdop,midway,above=1pt] {concat};
\node[mdch,above=1pt of KA-n] {4};
\draw[mdflow,rounded corners=4pt] (KA-e) -| (sum.south);
\node[rounded corners=3pt,draw=black!45,dashed,line width=0.6pt,inner sep=0pt,
      fit={($(K-w)+(-0.28,-1.30)$) ($(KS-e)+(0.18,2.08)$)}] (skipbox) {};
\node[mdch,anchor=north west] at ($(skipbox.north west)+(0.08,-0.06)$) {per band, $\times4$};
% ===================================================================== (b)
% same notation as (a); 192 -> 1.06, 32 -> .72, 25 -> .67
\def\yb{-6.62}
\def\hb{1.30}
\node[mdpanel] at (0,\yb+1.45) {(b)};
\mdfmap{bx}{1.40}{\yb}{0.72}{\hb}{mdLR}
\mdfmap{b1}{4.65}{\yb}{1.06}{\hb}{mdLR}
\mdfmap{b2}{8.05}{\yb}{0.67}{\hb}{mdLR}
\mdfmap{b3}{10.95}{\yb}{0.72}{\hb}{mdLR}
\node[mdsum] (bsum) at (13.05,\yb) {$+$};
\mdfmap{bo}{14.00}{\yb}{0.72}{\hb}{mdLR}
\draw[mdflow] (bx-e) -- (b1-w) node[mdop,midway,above=1pt] {conv $1\!\times\!1$\\ ReLU};
\draw[mdflow] (b1-e) -- (b2-w) node[mdop,midway,above=1pt] {conv $1\!\times\!1$};
\draw[mdflow] (b2-e) -- (b3-w) node[mdop,midway,above=1pt] {conv $3\!\times\!3$};
\draw[mdflow] (b3-e) -- (bsum);
\draw[mdflow] (bsum) -- (bo-w);
\draw[mdflow,rounded corners=5pt] ($(bx-s)+(0.06,0)$) -- ++(0,-0.48) -| (bsum.south)
  node[mdop,pos=0.25,above=1pt] {identity};
\node[mdch,above=1pt of bx-n] {32};
\node[mdch,above=1pt of b1-n] {192};
\node[mdch,above=1pt of b2-n] {25};
\node[mdch,above=1pt of b3-n] {32};
\node[mdch,above=1pt of bo-n] {32};
\node[rounded corners=4pt,draw=mdRes!85!black,line width=0.8pt,inner sep=0pt,
      fit={(0.62,\yb-1.20) (18.00,\yb+1.45)}] (bbox) {};
\node[font=\fontsize{9}{10.5}\selectfont\bfseries,text=mdRes!65!black,align=center,inner sep=0pt]
  at (16.75,\yb+0.12) {residual\\ block};
\end{tikzpicture}
\caption{The reconstruction network.
(a) The four-band input on the $0.1''$ grid is stretched with $\mathrm{arcsinh}$ at eight knees, giving $32$ channels.
The shared trunk (a $3\times3$ convolution, $16$ residual blocks, a second $3\times3$ convolution and a pixel shuffle) predicts a correction that is added to a per-band path, in which each band is upsampled from its own eight channels by a $5\times5$ convolution and a pixel shuffle.
A learned knee $k_b$ for each band converts the sum to electrons on the $0.05''$ grid.
Numbers above the blocks are channel counts; gray blocks lie on the input grid and green blocks on the output grid.
(b) One residual block: a $1\times1$ convolution widens $32$ channels to $192$ before the ReLU, a linear $1\times1$ convolution narrows them to $25$, and a $3\times3$ convolution returns $32$ channels, which are added to the block input.
The thumbnails show a Q1 lens candidate and its reconstruction (R\,=\,$H_{\rm E}$, G\,=\,mean of $Y_{\rm E}$ and $J_{\rm E}$, B\,=\,VIS).}
\label{fig:architecture}
\end{figure*}
